\section{Compressing repeated sums}
\label{sec:compression}

The first trial applies when random blocks have many repeated subset
sums. It retains only the distinct weights of seven disjoint blocks,
uses those small lists to generate four larger groups, and joins the
groups one residue class at a time. Because the domains are disjoint,
we can forget which subset produced a retained weight. A hard
capacity bounds memory; success requires only the residue class of
one fixed solution to fit.

This trial uses the classical stored-list and generated-stream
viewpoint of Schroeppel and Shamir~\cite{SS81}. Its modular
regeneration and overflow rule also belong to the dissection approach:
Dinur, Dunkelman, Keller, and Shamir~\cite{DDKS12} developed the
earlier framework under randomness assumptions, and Austrin, Kaski,
Koivisto, and M\"a\"att\"a~\cite{AKKM13} adapted dissection to
worst-case Subset Sum. The trial and its bounds here are proved
directly, including the dependence on the input bit length.

\subsection{Seven compacted blocks and four groups}

Take the first $7s$ indices of a uniform random permutation and divide
them into seven consecutive blocks $K_1,\ldots,K_7$ of size $s$.
For each $i$, enumerate and sort all its subset sums, then compact
repeated weights into a list $D_i$ of distinct actual weights.
Continue only if every list has at most $2^a$ entries, where
\begin{equation}
\label{comp:block-cap}
 a=\lfloor s-.0075n\rfloor.
\end{equation}
This preliminary enumeration stores at most $2^s$ weights at a time.
The seven retained lists are stored separately.

Of the remaining indices, give
\[
 r_R=\lfloor .05n\rfloor
 \quad\text{to the right and}\quad
 r_L=n-7s-r_R
 \quad\text{to the left}.
\]
Choose these domains in their permutation order. Split the right raw
domain into $V_{R1},V_{R2}$ of sizes $r_{R1},r_{R2}$ differing by
at most one, and call the left raw domain $V_L$.
All these sizes are nonnegative on sufficiently large main-branch
inputs.

For a domain $V$, write $\mathcal S(V)$ for its full list of indexed
subset-sum occurrences, retaining repetitions. For two lists $A,D$
write $A\oplus D$ for the list with one occurrence of weight $x+y$
for each pair $(x,y)\in A\times D$. Table~\ref{tab:compression}
specifies the four generated groups by their two stored lists.
A group stream enumerates the pair sums of the two lists in its row.
The eight stored lists use separate subdomains, and the four group
domains partition the input indices.

\begin{table}[htbp]
\centering
\begin{tabular}{@{}llll@{}}
\toprule
Group & Index domain & First stored list & Second stored list\\
\midrule
$\mathcal L_1$ & $K_5\cup K_6$ & $D_5$ & $D_6$\\
$\mathcal L_2$ & $K_7\cup V_L$ & $D_7$ & $\mathcal S(V_L)$\\
$\mathcal R_1$ & $K_1\cup K_2\cup V_{R1}$ & $D_1$
 & $D_2\oplus\mathcal S(V_{R1})$\\
$\mathcal R_2$ & $K_3\cup K_4\cup V_{R2}$ & $D_3$
 & $D_4\oplus\mathcal S(V_{R2})$\\
\bottomrule
\end{tabular}
\caption{The seven compacted block lists generate four groups. Each
row is streamed from its two stored lists. The raw and product lists
retain occurrences, while each $D_i$ retains only distinct weights.
No two rows share an input index.}
\label{tab:compression}
\end{table}

Sample $q$ uniformly from $\mathcal P_\ell$ with
$\ell=\lfloor .249n\rfloor$. For each integer
$z\in\{0,\ldots,q-1\}$, proceed as follows.
\begin{enumerate}
\item Generate $\mathcal L_1$ by its direct key modulo $q$ and
$\mathcal L_2$ by its complementary key $z-\text{weight}$
modulo $q$. Join equal keys using Lemma~\ref{lem:streams}, with
$p=1$ and $P=q$. Insert each resulting actual total weight into
a weight-only dictionary of capacity $C$.
\item If this dictionary overflows, discard it and move to the next
$z$. Otherwise retain its completed list of distinct weights.
An empty completed dictionary may also be skipped.
\item Generate and join $\mathcal R_1,\mathcal R_2$ in the same
way, with target residue $t-z\pmod q$. For every resulting actual
right weight $y$, search for $t-y$ in the left dictionary. A hit
returns YES. Otherwise free the dictionary and continue the residue
loop.
\end{enumerate}
A sampling failure also fails the trial. The bounded samplers and
operation counter are those of Section~\ref{sec:implementation}.

Every retained block weight has a subset witness in its block. The
four group domains are disjoint, so any choices witnessing an exact
left--right weight match combine into an indexed subset of weight
$t$. Thus the trial has one-sided correctness. It need not retain
witnesses or output them: its required output is the decision YES.

\subsection{Work, space, and success}

Let $\mu$ be the mean of $\defect(V)$ when $V$ is uniform among
the size-$s$ domains of the given input. The algorithm never computes
this mean. It runs this trial and the representation trial in the same
round; the following hypothesis identifies the case in which this
first trial supplies success.

\begin{lemma}[Compressed-block trial]
\label{lem:compressed}
For sufficiently large main-branch inputs, the compressed-block trial
uses at most $u^{15}2^{n/2}$ word operations on every random outcome,
including bounded sampling, and uses
$O(u(C+1)+B+2^s+2^{.14n})$ writable words. Every returned YES is
correct. If the input is a YES instance and
\begin{equation}
\label{alg:high-deficiency}
 \mu\ge .00765n,
\end{equation}
its success probability is bounded below by an absolute positive
constant.
\end{lemma}

\begin{proof}
First bound the stored and generated lists deterministically, assuming
the seven size tests pass. To retain the effects of floors, put
\[
 \delta=.12n-s\in[0,4),\qquad
 \epsilon=.05n-r_R\in[0,1).
\]
Then $r_L=.11n+7\delta+\epsilon$ and
$a\le.1125n-\delta$. Each stored product list on the right
has at most
\[
 2^{a+\lceil r_R/2\rceil}\le 2^{.1375n+1/2}
\]
occurrences. The left raw list has at most $2^{.11n+29}$
occurrences, and each $D_i$ has at most $2^{.1125n}$ entries.
Hence every stored sublist in Table~\ref{tab:compression} has size
at most $2^{.14n}$ for sufficiently large $n$.

Let $N_{L1},N_{L2},N_{R1},N_{R2}$ denote the lengths of the four
group streams, counting occurrences. Their direct product bounds are
\[
\begin{aligned}
 N_{L1}&\le2^{2a}\le2^{.225n},&
 N_{L2}&\le2^{a+r_L}\le2^{.2225n+25},\\
 N_{Ri}&\le2^{2a+r_{Ri}}\le2^{.25n+1/2}
 &&(i=1,2).
\end{aligned}
\]
More importantly, the products that count complete combinations obey
\begin{align}
 N_{R1}N_{R2}
 &\le2^{4a+r_R}
 \le2^{.5n-4\delta-\epsilon}\le2^{n/2},
 \label{comp:right-product}\\
 N_{L1}N_{L2}
 &\le2^{3a+r_L}
 \le2^{.4475n+4\delta+\epsilon}
 <2^{.4475n+17}.
 \label{comp:left-product}
\end{align}
These count all combinations, not just distinct resulting weights.

Each complete left combination has exactly one total residue modulo
$q$, and each right combination has exactly one complementary
residue $t-\text{weight}$. Across the whole residue loop, the
number of potential Cartesian matches is consequently at most
$N_{L1}N_{L2}+N_{R1}N_{R2}$. Overflow may stop a join partway, but
Lemma~\ref{lem:streams} still charges it to this same potential
count. Repeated weights therefore cause neither an uncounted replay
cost nor a loss of a combination.

Here is a concrete operation accounting. Generating an item, including
its weight, uses at most $O(n^2)$ operations; comparison sorting and
heap operations use $O(n)$ per item or advance; trie operations use
$O(u)$ per key. The stream states for a group have $O(2^{.14n})$
words and are therefore $O(B)$ for large $n$. Thus the match-dependent
term in~\eqref{eq:tie-cost} has $O(n)$ overhead per potential match;
its buffer and stream-length terms are charged separately below.
We may even initialize $O(u(C+1))$ trie space afresh at each residue.
If $A$ is the total number of stored sublist entries and
$N=N_{L1}+N_{L2}+N_{R1}+N_{R2}$, all nonsampling work is bounded by
\begin{equation}
\label{comp:operation-ledger}
 O\left(u^3\left[
  2^s+A+q\bigl(N+B+u(C+1)+1\bigr)
  +N_{L1}N_{L2}+N_{R1}N_{R2}\right]\right).
\end{equation}
The $2^s$ term pays the seven initial enumerations and compactions,
the terms multiplied by $q$ pay regeneration and initialization, and
the final two terms pay every potential match and dictionary action.
The bound also covers a failed size test.

Since $q<2^{.249n}$, the largest regenerated-stream exponent is
$.249+.25=.499$. The buffer and dictionary initialization exponents
are $.439$ and $.448$, respectively. Equations
\eqref{comp:right-product}--\eqref{comp:left-product} bound the
remaining products, so~\eqref{comp:operation-ledger} is
$O(u^4 2^{n/2})$. A permutation uses at most $O(n^4)$ bounded
range attempts and elementary operations, and the one prime search
costs at most $\poly(n)2^{.129n}$ with the fixed elementary
trial-division implementation of Lemma~\ref{lem:samplers}.
In particular it is $O(n^4 2^{.129n})$. These smaller terms leave
the total below $u^{15}2^{n/2}$ for sufficiently large $n$.
This deterministic bound is below the trial cap; it does not assume
any favorable bucket sizes or random events.

For space, the initial enumerations use $O(2^s)$ words, the stored
sublists and stream states use $O(2^{.14n})$, tie replay uses $O(B)$,
and the one retained dictionary uses $O(u(C+1))$. Sorting needs only
linear auxiliary arrays. The permutation and loop state are absorbed
by these bounds. In particular, overflow enforces the space bound on
every outcome.

It remains to show that one solution bucket fits in the high-mean
case. Each random $K_i$ is marginally a uniform size-$s$ domain;
the seven blocks need not be independent. Lemma~\ref{lem:deficiency},
Chebyshev's inequality, and a union bound give
\[
 \Pr\bigl[\defect(K_i)\ge .00755n
             \text{ for all }i\bigr]=1-O(1/n)
\]
under~\eqref{alg:high-deficiency}. These inequalities imply
$N(K_i)\le2^{s-.00755n}\le2^a$ for sufficiently large $n$,
so all tests pass, including their floors.

Fix an original solution before sampling the partition. Conditional
on any passing partition, let $x$ be its left weight. There are fewer
than $2^{.4475n+17}$ distinct left weights in all. For each different
one $x'$, Lemma~\ref{lem:prime} and
$\ell=\lfloor .249n\rfloor$ give
\[
 \Pr[x'\equiv x\pmod q]
 \le 8u^2 2^{-\ell}\le u^3 2^{-.249n}
\]
for large $n$. By linearity of expectation, the number $Y$ of
other distinct weights in $x$'s residue class satisfies
\[
 \mathbb E Y\le u^3 2^{.1985n+17}.
\]
Because $C\ge2^{.199n-1}$ and $\log u\le n/10^9$,
\[
 \frac{\mathbb E Y}{C}
 \le 2^{(3/10^9-.0005)n+18}=o(1).
\]
Markov's inequality implies $Y+1\le C$ with probability tending
to one. The complete dictionary in that residue then retains $x$,
and the right solution weight produces a hit. Earlier residues can
only produce an earlier correct YES or be skipped.

This proves success with probability $1-o(1)$ in the ideal experiment.
The permutation and prime draws have negligible total bounded-sampler
coupling loss by Lemma~\ref{lem:samplers}; the deterministic work
estimate prevents the operation counter from aborting a completed
trial for sufficiently large $n$. Thus an absolute positive success
probability remains for the implemented trial.
\end{proof}

Compression has now handled the case in which random blocks have few
distinct sums. When $\mu<.00765n$, the next trial instead uses the
many distinct sums to preserve a representation of a fixed solution
through several modular filters.
